\scanpage{5}
\uline{Cor. 4} Les tores déployés \inserted{de dim $n$} de $\underline{\mathrm{Cr}}_{n,k}$ sont conjugués.

\uline{Cor.} 5 Ils sont maximaux, ils sont leurs propres \unsure{centralisateur|centralisateurs}, et le normalisateur est un groupe \deleted{\unsure{lisse|tore}} lisse, et $N/T\simeq\underline{\mathrm{Aut}}_{\mathrm{gr}}(T)$.

\uline{Cor.} 6 $T$ est un tore maximal de $\underline{\mathrm{Cr}}_{n,k}$

\uline{Proposition} $T$ déployé, $X$ irréductible, pseudo-opération fidèle. Alors $X\overset{\text{\unsure{pseudo}}}{\sim}T\times Y$, où $Y$ est un schéma irréductible (unique à pseudo-isom près).

\uline{Cor} 1 Soit $0\leq d\leq n$. Conditions équivalentes

a) \unsure{Tous|Tout|Im} $G_m^d\subset\underline{\mathrm{Cr}}_{n,k}$ est contenu dans un $G_m^n$

b) Toute extension $K/k$ de degr. tr. $d$, \unsure{dt|et} $K[t_1,-,t_{n-d}]$ pure, est pure (``Lüroth'').

[\uline{NB} \quad OK si $d=0$, $d=1$, $d=2$ $\left|\begin{array}{l}\text{car. nulle}\\ k\ \text{alg. clos}\end{array}\right.$ (Castelnuovo).

\hspace*{4em}$d=n$

Marche si $n\leq 3$ \quad ]

\uline{Pb} Classifier les $G$ alg. lisses $\subset\underset{\sim}{\mathrm{Cr}}_{n,k}$ qui contiennent un tore déployé de dim $n$.

\uline{\unsure{Th. des actions|Th. de dém.|Th. de structure}}

On va fixer $T$, et chercher
\[
T\subset G\subset\underset{\sim}{\mathrm{Psaut}}(T)
\]
L'opération de $G$ sera transitive. On veut espaces homogènes $X$ sous $G$, de dim $n$, de telle façon que $T$ opère fidèlement avec une orbite ouverte

% archive sheet number: 4
